\documentclass[12pt]{article}
\usepackage[T1]{fontenc}
\usepackage[utf8]{inputenc}
\usepackage{lmodern}
\usepackage{textcomp}
\usepackage{enumitem}
\usepackage{geometry}
\usepackage{graphicx}
\usepackage{amsmath, amssymb}
\geometry{margin=1in}
\begin{document}
\begin{center}
Philosophy - Undergraduate Level Assessment 17 \
Total Marks: 40 \
Answer all questions.
\end{center}

\section*{Multiple Choice (10 marks)}
\begin{enumerate}[label=\arabic*.]
\item Which leader is known for spreading the dharma of non-violence?
\begin{enumerate}[label=(\alph*)]
    \item Ashoka
    \item Nelson Mandela
    \item Martin Luther King Jr.
    \item Ngo Dinh Diem
    \item Adolf Hitler
\end{enumerate}
\item According to Jaina traditions, what does the term ajiva mean?
\begin{enumerate}[label=(\alph*)]
    \item Sound
    \item Non-living
    \item Non-matter
    \item Non-soul
    \item Consciousness
\end{enumerate}
\item When someone responds to your argument with a sarcastic statement such as, "Yeah, right. Like that's ever going to happen," that person may have committed which fallacy?
\begin{enumerate}[label=(\alph*)]
    \item tu quoque
    \item appeal to ignorance
    \item argumentum ad populum
    \item false cause
    \item appeal to indignation
\end{enumerate}
\item Hobbes defines injustice as:
\begin{enumerate}[label=(\alph*)]
    \item not adhering to societal norms.
    \item being dishonest in any situation.
    \item failure to perform one's covenant.
    \item treating another person as a mere means.
    \item manipulating others for personal gain.
\end{enumerate}
\item According to Mill, to determine whether one pleasure is more valuable than another, we must \_\_\_\_\_.
\begin{enumerate}[label=(\alph*)]
    \item determine which pleasure most experienced people prefer
    \item consult science
    \item consult religious leaders
    \item determine which one is objectively most pleasurable
    \item measure the intensity of each pleasure
\end{enumerate}
\end{enumerate}

\section*{True / False (10 marks)}
\textit{Answer True or False}
\begin{enumerate}[label=\arabic*.]
\setcounter{enumi}{5}
\item True or False: Nagel posits that absolutism holds we should only sometimes prevent murder based on contextual factors.
\item True or False: Augustine asserts that Academic skepticism is false and can be effectively refuted through rational argumentation.
\item True or False: Socrates believed that an unexamined life primarily harms the individual's reputation in society.
\item True or False: Aristotle asserts that both virtue and vice are within our power and subject to our control.
\item True or False: According to Hume, morality is ultimately based on cultural norms that shape societal values and behaviors.
\end{enumerate}

\section*{Long Form Response (20 marks)}
\textit{Answer each question with a clear and complete explanation.}
\begin{enumerate}[label=\arabic*.]
\setcounter{enumi}{10}
\item Discuss the concept of ziran as described by Zhuangzi, focusing on its implications for compassion in philosophical discourse. How does this understanding shape human relationships and ethical behavior?
\item Discuss the philosophical principles underlying Ikebana, the Japanese art of flower arranging, and analyze how it reflects broader concepts of beauty and nature in Japanese culture.
\end{enumerate}

\vfill
\noindent\textit{End of Paper}
\end{document}